%% kapa: exam=1 round=25 year=2014 period=2 form=A subject=accounting subject_ko=회계학 seq=38 number=78
%% source: Q-Net 감정평가사 기출문제 (cst00309) artlSeq=5101207 file=extracted/1008/5101207/2014년 제25회 감정평가사 1차시험 기출문제/2교시(A형).hwp
%% conversion: kordoc 4.24.51 -> tools/gen1.py (auto); review=proofread-sonnet; flags=none
\begin{question}{78}
(주)감평은 20×1년 초에 신제품을 개발하여 최초 10단위를 생산하였으며, 최초 10단위의 신제품 생산과 관련된 제조원가는 다음과 같다.
\par{\small\centering
\begin{tabular}{@{}l@{\hspace{3em}}r@{}}
\hline
직접재료원가 & ￦8,000 \\
직접노무원가(시간당 ￦10) & 3,000 \\
변동제조간접원가 (직접노무시간에 비례하여 발생) & 9,000 \\
고정제조간접원가 배부액 & 2,000 \\
\hline
\end{tabular}\par}
신제품 생산에는 90\% 누적평균시간 학습효과가 있는 것으로 분석되었다. 추가로 30단위를 생산하여 총 40단위를 생산할 경우, 총 40단위에 대해 예상되는 변동원가는? (단, (주)감평은 신제품을 생산할 수 있는 충분한 여유설비를 확보하고 있다.) \par
\begin{choices}
  \item ￦70,880
  \item ￦72,880
  \item ￦76,680
  \item ￦78,680
  \item ￦80,880
\end{choices}
\end{question}
\answer{1}
